\documentclass[aspectratio=169]{beamer}
\usepackage{hyperref}
\usepackage[final]{listings}
\usepackage{color}

% Define colors for listings
\definecolor{codegreen}{rgb}{0,0.6,0}
\definecolor{codegray}{rgb}{0.5,0.5,0.5}
\definecolor{codepurple}{rgb}{0.58,0,0.82}
\definecolor{backcolour}{rgb}{0.95,0.95,0.92}

\lstdefinestyle{mystyle}{
    backgroundcolor=\color{backcolour},   
    commentstyle=\color{codegreen},
    keywordstyle=\color{magenta},
    numberstyle=\tiny\color{codegray},
    stringstyle=\color{codepurple},
    basicstyle=\ttfamily\footnotesize,
    breakatwhitespace=false,         
    breaklines=true,                 
    captionpos=b,                    
    keepspaces=true,                 
    numbers=left,                    
    numbersep=5pt,                  
    showspaces=false,                
    showstringspaces=false,
    showtabs=false,                  
    tabsize=2
}

\lstset{style=mystyle}

\title{An In-Depth Guide to ChemPy: Python's Chemistry Library}
\author{}
\date{}

\begin{document}

\frame{\titlepage}

\begin{frame}
    \frametitle{Table of Contents}
    \tableofcontents
\end{frame}

\section{Introduction to ChemPy}

\begin{frame}
    \frametitle{What is ChemPy?}
    ChemPy is an open-source Python library for chemistry computations. It allows users to simulate chemical kinetics, equilibrium, and other processes using a flexible and intuitive interface. ChemPy is built on top of Numpy and Scipy, ensuring high performance and integration with other scientific libraries in Python.
\end{frame}

\begin{frame}[fragile]
    \frametitle{Installation}
    To install ChemPy, you can use pip:
    \begin{lstlisting}[language=bash]
pip install chempy
    \end{lstlisting}
\end{frame}

\section{Core Components}

\begin{frame}[fragile]
    \frametitle{Chemical Reactions}
    ChemPy provides tools to define and manipulate chemical reactions. You can represent reactions in a human-readable format and then use these representations to perform various calculations.
\end{frame}

\begin{frame}[fragile]
    \frametitle{Example: Defining a Simple Reaction}
    \begin{lstlisting}[language=python]
from chempy import Reaction

# Define the reaction: 2 H2 + O2 -> 2 H2O
reaction = Reaction.from_string("2 H2 + O2 -> 2 H2O")
print(reaction)
    \end{lstlisting}
\end{frame}

\begin{frame}
    \frametitle{Equilibrium Calculations}
    ChemPy can perform equilibrium calculations to find the concentrations of species at equilibrium for a given reaction.
\end{frame}

\begin{frame}[fragile]
    \frametitle{Example: Equilibrium Concentration}
    \begin{lstlisting}[language=python]
from chempy.equilibria import EqSystem
from chempy.thermodynamics import GibbsEqConsts
from chempy.util.units import to_unitless, default_units as u

# Define the reactions and equilibrium constants
reactions = [
    Reaction.from_string("H2 + I2 -> 2 HI")
]
eq_consts = GibbsEqConsts([1e-2])

# Create an equilibrium system
eqsys = EqSystem(reactions, eq_consts)

# Initial concentrations
initial_concs = {'H2': 1 * u.molar, 'I2': 1 * u.molar, 'HI': 0 * u.molar}

# Solve the equilibrium concentrations
eq_concs = eqsys.root(initial_concs)
print(to_unitless(eq_concs, u.molar))
    \end{lstlisting}
\end{frame}

\begin{frame}[fragile]
    \frametitle{Kinetic Simulations}
    ChemPy allows for the simulation of reaction kinetics using ordinary differential equations (ODEs). This is useful for studying the time evolution of reaction systems.
\end{frame}

\begin{frame}[fragile]
    \frametitle{Example: Reaction Kinetics}
    \begin{lstlisting}[language=python]
from chempy.kinetics.ode import get_odesys
from chempy.kinetics.rates import ArrheniusParam
from chempy.kinetics import ReactionRates

# Define the reaction and rate
reaction = Reaction.from_string("2 H2 + O2 -> 2 H2O")
rate = ArrheniusParam(A=1e10, Ea=50e3)

# Create the ODE system
odesys = get_odesys([reaction], [rate])

# Initial concentrations and time points
initial_concs = {'H2': 2.0, 'O2': 1.0, 'H2O': 0.0}
time_points = [0, 1, 2, 3, 4, 5]

# Solve the ODEs
result = odesys.integrate(time_points, initial_concs)
print(result)
    \end{lstlisting}
\end{frame}

\section{Advanced Features}

\begin{frame}
    \frametitle{Thermodynamics}
    ChemPy includes tools for thermodynamic calculations, such as Gibbs free energy, enthalpy, and entropy changes for reactions.
\end{frame}

\begin{frame}[fragile]
    \frametitle{Example: Gibbs Free Energy}
    \begin{lstlisting}[language=python]
from chempy.thermodynamics import Gibbs

# Define the reaction
reaction = Reaction.from_string("2 H2 + O2 -> 2 H2O")

# Gibbs free energy change
dG = Gibbs([reaction])
print(dG(298.15))  # at 298.15 K
    \end{lstlisting}
\end{frame}

\begin{frame}[fragile]
    \frametitle{Electrochemistry}
    ChemPy also provides functionality for electrochemical calculations, including cell potentials and Nernst equations.
\end{frame}

\begin{frame}[fragile]
    \frametitle{Example: Electrochemical Cell}
    \begin{lstlisting}[language=python]
from chempy.electrochemistry import ElectrochemicalCell
from chempy.thermodynamics import nernst

# Define half-cell reactions
half_cells = [
    Reaction.from_string("Zn -> Zn2+ + 2 e-"),
    Reaction.from_string("Cu2+ + 2 e- -> Cu")
]

# Standard electrode potentials
E0 = [0.76, 0.34]  # Zn/Zn2+ and Cu/Cu2+ standard potentials

# Create the electrochemical cell
cell = ElectrochemicalCell(half_cells, E0)

# Nernst equation at standard conditions
cell_potential = cell.potential(298.15)
print(cell_potential)
    \end{lstlisting}
\end{frame}

\section{Fun Examples}

\begin{frame}[fragile]
    \frametitle{Simple Reaction Kinetics}
    Let's simulate the kinetics of the reaction between hydrogen and oxygen to form water:
    \begin{lstlisting}[language=python]
# Already defined earlier
time_points = [0, 1, 2, 3, 4, 5]
result = odesys.integrate(time_points, initial_concs)

# Plot the results (requires matplotlib)
import matplotlib.pyplot as plt

plt.plot(result.tout, result.yout[:, 0], label='[H2]')
plt.plot(result.tout, result.yout[:, 1], label='[O2]')
plt.plot(result.tout, result.yout[:, 2], label='[H2O]')
plt.xlabel('Time [s]')
plt.ylabel('Concentration [M]')
plt.legend()
plt.show()
    \end{lstlisting}
\end{frame}

\begin{frame}[fragile]
    \frametitle{Buffer Solutions}
    Let's simulate the pH of a buffer solution made of acetic acid and sodium acetate:
    \begin{lstlisting}[language=python]
from chempy.kinetics.ode import integrate_ode
from chempy.kinetics.rates import MassAction
from chempy.equilibria import EqSystem

# Define the buffer reaction: CH3COOH <=> CH3COO- + H+
buffer_reaction = Reaction.from_string("CH3COOH <=> CH3COO- + H+")
buffer_eq_const = 1.8e-5  # Acetic acid dissociation constant

# Create equilibrium system
buffer_eqsys = EqSystem([buffer_reaction], [buffer_eq_const])

# Initial concentrations
initial_concs = {'CH3COOH': 0.1, 'CH3COO-': 0.1, 'H+': 1e-7}

# Solve the equilibrium concentrations
buffer_eq_concs = buffer_eqsys.root(initial_concs)
print(buffer_eq_concs)
    \end{lstlisting}
\end{frame}

\begin{frame}[fragile]
    \frametitle{Electrochemical Cell Simulation}
    Simulate the potential of a Daniell cell (Zn/Cu):
    \begin{lstlisting}[language=python]
# Already defined earlier
cell_potential = cell.potential(298.15)
print(f"Cell potential at 298.15 K: {cell_potential} V")
    \end{lstlisting}
\end{frame}

\section{Conclusion}

\begin{frame}
    \frametitle{Conclusion}
    ChemPy is a versatile library for chemistry-related computations in Python. From simple reaction kinetics to complex electrochemical simulations, ChemPy provides a comprehensive toolkit for chemists and educators. By integrating with other scientific libraries in Python, it offers high performance and flexibility for a wide range of applications.
\end{frame}

\end{document}
